\subsection{ORTHOGONAL REFLECTIONS IN A EUCLIDEAN AFFINE SPACE}

We retain the notation of the preceding number, and let E be an affine space of which V is the space of translations. Giving the form B on V provides E with the structure of a euclidean space (Algebra, Chap. IX, §6, no. 6).

Let H be a non-isotropic hyperplane of E. The symmetry with respect to H (Algebra, Chap. IX, §6, no. 6) is also called the orthogonal reflection with respect to H; we often denote it by \(s_{H}\) . We have \(s_{H}^{2}=1\) and \(s_{H}\) is the unique displacement (loc. cit., Def. 3) of E, distinct from the identity and leaving fixed the elements of H. The automorphism of V associated to \(s_{H}\) is the orthogonal reflection with respect to the direction of H (which is a non-isotropic hyperplane of V).

Every \(x\) in \(\mathbf{E}\) can be written uniquely in the form \(x = h + v\) , with \(h \in \mathbf{H}\) and \(v \in \mathbf{V}\) orthogonal to \(\mathbf{H}\) ; we have

\[
s _ {\mathrm{H}} (h + v) = h - v.
\]

PROPOSITION 5. Let H and H' be two parallel, non-isotropic hyperplanes of E. There exists a unique vector \(v \in V\) orthogonal to H and such that \(H' = H + v\) . The displacement \(s_{H'} s_{H}\) is the translation by the vector 2v.

The existence and uniqueness of v are immediate. The automorphism of V associated to \(s_{H'}s_{H}\) is the identity; thus \(s_{H'}s_{H}\) is a translation. On the other hand, let \(a \in H'\) ; then \(a - v \in H\) and

\[
s _ {\mathrm {H^ {\prime}}} s _ {\mathrm{H}} (a - v) = s _ {\mathrm {H^ {\prime}}} (a - v) = a + v = (a - v) + 2 v,
\]

showing that \(s_{H'}s_{H}\) is the translation by the vector 2v.

COROLLARY. Let H and H' be two distinct, parallel, non-isotropic hyperplanes. If K is of characteristic zero (resp. p \textgreater{} 0, with \(p \neq 2\) ), the group of displacements of E generated by \(s_{H}\) and \(s_{H'}\) is an infinite dihedral group (resp. a dihedral group of order 2p).

Indeed, by Prop. 2 of Chap. IV, §1, no. 2, it suffices to show that \(s_{\mathrm{H}'}s_{\mathrm{H}}\) is of infinite order (resp. order \(2p\) ), which is clear.

Remark. We retain the notation of Prop. 5 and assume in addition that K = R. Put \(s = s_{H}\) and \(s' = s_{H'}\) . Let \(H_{n}\) be the hyperplane \(H + n.v\) and let \(C_{n}\) be the set of points of E of the form \(a + \xi.v\) with \(a \in H\) and \(n < \xi < n + 1\) . The \(C_{n}\) are connected open sets forming a partition of \(E - \bigcup_{n} H_{n}\) . They are therefore the chambers defined by the system \(\mathfrak{H} = (\mathrm{H}_{n})_{n \in \mathbf{Z}}\) in E. The translation \((s's)^{n}\) transforms the chamber \(C = C_{0}\) into the chamber \(C_{2n}\) , and since \(s(\mathrm{C}_{0}) = \mathrm{C}_{-1}\) , we have \((s's)^{n}s(\mathrm{C}) = \mathrm{C}_{2n-1}\) . It follows that the dihedral group W generated by s and \(s'\) permutes the chambers \(C_{n}\) simply-transitively. Moreover, as we shall now show, if the chambers C and w(C) are on opposite sides of H (for \(w \in W\) ), we have \(l(sw) = l(w) - 1\) (the lengths being taken with respect to S = \{s, s'\} (Chap. IV, §1, no. 1)). Indeed, we then have \(w(\mathrm{C}) = C_{n}\) for some n \textless{} 0. If n = -2k, then \(w = (ss')^{k}\) and \(sw = (s's)^{k-1}s'\) , so \(l(w) = 2k\) and \(l(sw) = 2k - 1\) (Chap. IV, §1, no. 2, Remark). If n = -2k - 1, then \(w = (ss')^{k}s\) and \(sw = (s's)^{k}\) , so \(l(w) = 2k + 1\) and \(l(sw) = 2k\) .

